\documentstyle[%


righttag%


,11pt%


,amssymb%


,russian%               remove in Eng.


,izvru%                 Set izven% in Eng.


]{amsart}





%





\newcommand{\tts}[1]{}


\numberwithin{equation}{section}





\theoremstyle{definition}


\theorembodyfont{\rm}


\newtheorem{notes}{\hskip\parindent Замечание}[section]


\newtheorem{defns}{\hskip\parindent Определение}[section]





%\hoffset=-15mm%                  For Eng.


\setcounter{page}{10}


\god{2001}


\nomer{12~(475)} %                        Remove in Eng.


%\nomer{Vol.\,45, No.\,12}%               Set in Eng.


%\str{\pageref{begin}--\pageref{end}}%  Set in Eng.


\udk{517.852}





\author{\it Р.\,Габасов, Ф.М.\,Кириллова}


% NB:   ^^ remove in Eng.


\title{К проблеме синтеза оптимальных систем}


\thanks{Работа выполнена при финансовой поддержке Белорусского фонда


фундаментальных исследований (проект \No\,Ф99Р-002).}





\date{}


%\pagestyle{myheadings}%         Set in Eng.


\pagestyle{plain} %              Remove in Eng.


\begin{document}


%\thispagestyle{plain}%          Set in Eng.


\maketitle


\label{begin}








В теории управления проблема синтеза оптимальных систем (оптимальное


управление по принципу обратной связи) является одной из основных.


Конструктивное решение этой проблемы позволяет


не только эффективно решать прикладные задачи оптимизации


динамических систем, но и исследовать ряд задач, которые по своей


постановке не являются экстремальными (стабилизация,


демпфирование, амортизация и др.). Согласно ([1], гл.\,1, с.\,16) при


управлении в условиях постоянно действующих возмущений


используются три принципа управления: 1)~принцип обратной связи,


2)~прямой связи, 3)~прямо-обратной связи


(комбинированный принцип). Принцип~1), использующий при формировании


управлений выходные сигналы системы, является более универсальным,


чем~2). Однако существуют ситуации,


когда целесообразно применять второй принцип, основанный на


доступных измерениях возмущений (входных сигналов).


Комбинированный принцип управления базируется на~1),~2)


и реализуется в виде прямых и обратных связей (использует доступные измерения


входных и выходных сигналов).





В данной статье для линейных нестационарных динамических систем


излагается подход к проблеме синтеза оптимальных систем,


функционирующих в условиях постоянно действующих возмущений,


который базируется на процедуре коррекций оптимальных программных


управлений и быстрых алгоритмах оптимизации. В классе дискретных


управлений обоснованы алгоритмы построения реализаций оптимальных


дискретных обратных связей, прямых и


комбинированных связей. В \S\,1 излагается подход к построению


оптимальных замыкаемых ([2], гл.\,1, с.\,18; [3]) и замкнутых обратных


связей, обеспечивающих оптимальные гарантированные результаты для


систем управления, подверженных действию ограниченных возмущений


нестохастической природы. В \S\,2 изучаются оптимальные размыкаемые


связи, строящиеся по результатам измерений возмущений (принцип


прямой связи). В \S\,3 при построении реализации оптимальных связей


применяется комбинированный принцип управления, использующий


входные и выходные сигналы.








\section{Синтез оптимальных обратных связей для динамических


систем\protect \\


с возмущениями}





Пусть $T=[t_{*}, t^{*}]$ --- промежуток


управления, $h=(t^{*}-t_{*})/N$ --- период квантования времени, $N$


--- натуральное число, $T_{h}=\{t_{*},\ t_{*}+h,\ t_{*}+(N-1)h\}$.


Функцию $u(t)$, $t \in T$, назовем дискретным управлением, если


$$


u(t)=u(t_{*}+kh),\ \; t \in


[t_{*}+kh,t_{*}+(k+1)h[,\ \; k=\overline{0,N-1}.


$$





В классе дискретных управлений рассмотрим задачу оптимального


управления


\begin{gather}


c'x(t^{*})\rightarrow \max,\quad \Dot x=A(t)x+b(t)u+d(t)w,\quad


x(t_*)=x_{0},\\


%


x(t^*) \in X^*=\{x \in R^n:g_*\le Hx \le g^*\},\quad u(t) \in U=\{u \in


R:|u| \le 1\}, \notag \\


%


w(t) \in W=\{w \in R:|w| \le 1\},\quad t \in T,\notag


\end{gather}


где $x \in R^n$, $u\in R$, $H \in R^{m \times n}$, $A(t)$, $b(t)$, $d(t)$,


$t \in T$, --- кусочно-непрерывные функции, $w=w(t)$ --- неизвестное


возмущение, $w \in R$, $t \in T$. Пусть до начала процесса


управления известно, что 1)~в процессе управления неизвестное


возмущение  может реализоваться в виде любой кусочно-непрерывной


функции $w(t) \in W$, $t \in T$; 2)~в заданные моменты замыкания


\begin{equation}


T^p=\{t^i\in T_h,\ i=\overline{1,p}\},\quad t_* < t^1<\dotsb<t^p<t^*,


\end{equation}


будут доступны измерению состояния $x(t)$, $t \in T^p$, системы;


3)~в моменты времени  $\tau \in T_{h} \setminus T^p$ текущие


состояния $x(\tau)$ могут быть и доступными и недоступными.





При этих условиях неформальная постановка задачи состоит в


построении программных и типа обратной связи (позиционных)


управлений, под действием которых динамическая система, во-первых, в


момент $t^*$ достигает множества $X^*$  независимо от


реализовавшегося возмущения, а во-вторых, доставляет критерию качества


$c'x(t^*)$ максимальное гарантированное значение. Для строгой


постановки задачи проведем некоторые построения. По совокупности


(1.2) построим множества


\begin{equation}


X^i,\quad i=\overline{0,p+1}.


\end{equation}





Положим $X^{p+1}=X^*$.  Множество $X^i$  построим по множеству


$X^{i+1}:z \in X^i$ тогда и только тогда, когда существует


доступное управление $u(t) \in U$, $t \in [t^i, t^{i+1}[$,


переводящее в момент $t^{i+1}$ динамическую систему из состояния


$x(t^i)=z$ на множество $X^{i+1}$ при любом возмущении $w(t) \in


W$, $t \in [t^i, t^{i+1}]$.





Ясно, что задача (1.1) имеет решение в том и только том случае,


если $x_0 \in X^0$.





\begin{notes} Множества (1.3) можно определить и с помощью


детерминированной модели


\begin{equation}


\Dot x =A(t)x+b(t)u.


\end{equation}


\end{notes}





Существуют [3], [4] такие сужения


\begin{equation}


X_0^i,\quad i=\overline{1,p+1},


\end{equation}


множеств (1.3), что $z \in X^i$ тогда и только тогда, когда


найдется доступное управление $u(t) \in U$, $t \in [t^i,t^{i+1}]$,


которое переводит в момент $t^{i+1}$ систему (1.4) из состояния


$x(t^i)=z$ в состояние $x(t^{i+1}) \in X^{i+1}_0$.





Доступное управление $u(t) \in U$, $t \in T$, назовем


допустимым программным управлением задачи


(1.1), если соответствующая ему траектория


$x_0(t)$, $t \in T$, системы (1.4), $x_0(t_*)=x_0$, в момент


$t^*$ попадает на множество $X_0^*$.





По построению допустимое управление переводит систему (1.1) в


момент $t^*$ на множество $X^*$ при любых возможных реализациях


возмущения. Чтобы оценить качество допустимых управлений, введем


новую совокупность множеств


\begin{equation}


X^{i,\alpha},\quad i =\overline{0,p+1}.


\end{equation}


Положим $X^{p+1,\alpha}=X^* \cap \{x \in R^n:c'x \ge \alpha\}$.


Остальные множества из (1.6) построим по правилам построения


множеств (1.3). Пусть


\begin{equation}


X^{i,\alpha}_0,\quad i =\overline{1,p+1},


\end{equation}


--- сужения множеств (1.6). При достаточно малых $\alpha$ множества


(1.7) совпадают с (1.5). Обозначим через $\alpha (t^i)$


максимальное число $\alpha$, при котором $x_0(t^i) \in


X^{i,\alpha}_{0}$.





\begin{defns}


Число


$$


\alpha(u)=\min\alpha(t^i),\quad i = \overline{1,p+1},


$$


называется (гарантированным) значением критерия качества на


допустимом управлении $u(t)$, $t \in T$.


\end{defns}





\begin{defns}


Допустимое управление $u^0(t)$, $t \in T$, называется


оптимальным программным управлением задачи (1.1), если


$$


\alpha^0=\alpha(u^0)=\max\alpha(u),


$$


где максимум берется по всем допустимым управлениям.


\end{defns}





Пусть $\tau$, $t^i \le \tau < t^{i+1}$, ---


текущий момент времени. Обозначим через


$\underline{\tau}=\underline{\tau}(\tau)$ ближайший предыдущий к


$\tau$ момент из $t^i \le \underline{\tau} \le \tau$, в котором


было доступно измерению состояние $x(\underline{\tau})$;


$u^0(t\mid\underline{\tau},z)$, $t \in [\underline{\tau},t ^*]$, ---


оптимальное программное управление в задаче


\begin{gather}


c'x(t^{*})\rightarrow \max,\quad \Dot x=A(t)x+b(t)u+d(t)w,


\quad x(\tau)=z,\\


%


x(t^*) \in X^*,\ \; u(t) \in U,


\ \; w(t) \in W,\ \; t \in [\underline{\tau}, t^*],\notag


\end{gather}


$X_{\tau}$ --- множество всех $z \in R^n$, для которых задача


(1.8) имеет решение.





\begin{defns}


Функцию


\begin{equation}


u^0(\tau,z)=u^0(\tau \mid \underline{\tau}(\tau),z),\ \; z \in


X_{\underline{\tau}},\ \; \tau \in T_{h},


\end{equation}


будем называть оптимальным управлением типа замыкаемой дискретной


обратной связи.


\end{defns}





Замкнем систему (1.1) обратной связью (1.9). Под траекторией


замкнутой системы


$$


\Dot x =A(t)x+b(t)u^0(t,x(\underline{\tau}(t)))+d(t)w,\quad x(t_*)=x_0,


$$


будем понимать решение линейного уравнения


$$


\Dot x =A(t)x+b(t)u^*(t)+d(t)w,\quad x(t_*)=x_0,


$$


с управлением


\begin{equation}


u^*(t)=u^0(t_*+kh,x(\underline{\tau}(t_*+kh))),\ \; t \in


[t_*+kh,t_*+(k+1)h[,\ \; k=\overline{0,N-1}.


\end{equation}





Функцию (1.10) назовем {\it реализацией} обратной связи (1.9) в


конкретном процессе управления. Если в каждый текущий момент $\tau


\in T_h$ время на вычисление $u^*(\tau)$ не превосходит $h$, то


будем говорить, что обратная связь (1.9) реализуется в режиме


реального времени. Устройство, способное выполнять такую работу,


называется оптимальным регулятором. Таким образом, задача


оптимального синтеза свелась к описанию алгоритма работы


оптимального регулятора. Предлагаемый ниже алгоритм основан на


двух процедурах: 1)~решении вспомогательных линейных задач


оптимального управления; 2)~коррекции (доводки) приближенного


решения.


\medskip





1.1. {\it Решение вспомогательных задач}. Рассмотрим одно из


множеств $X^*_{\gamma}$ семейства


$$


X_\gamma^*=\{x \in R^n:g_{*l}-\gamma \le H'_{(l)}x \le


g_l^*+\gamma,\ l=\overline{1,m}\},\ \text{$H_{(l)}$ --- $l$-я строка


матрицы $H$}.


$$


Начальное значение $\gamma^* \ge 0$ выберем небольшим. Пусть


$\alpha \in [\alpha_*, \alpha^*]$, $\alpha_*=\min c'x$, $x \in X^*$;


$\alpha^*=\max c'x$, $x \in X^*$. Определим множество


$X^{*,\alpha}_{\gamma}=X^*_{\gamma} \cap \{x \in R^n: c'x \ge


\alpha\}$.





Введем конечную совокупность единичных векторов


\begin{equation}


f_j,\quad j = \overline{1,q},


\end{equation}


и рассмотрим промежуток $[t^p,t^*]$. Для каждого вектора из (1.11)


решим задачу гарантированной оптимизации


\begin{gather}


\max_{z,u}f'_{j}z=f'_{j}x_{j}^{p,\alpha}=\beta_{j}^{\alpha}(t^p),\notag\\


%


\Dot x =A(t)x+b(t)u+d(t)w,\quad x(t^p)=z,\\


%


x(t^*)\in X_{\gamma^*}^{*,\alpha},\ \; u(t) \in U,\ \; w(t) \in W,


\ \; t \in [t^p, t^*].\notag


\end{gather}


Согласно [3] по элементам задачи (1.12) можно подсчитать такие


числа $\alpha_0$, $\underline{\alpha}_{\,l}^{*},


\overline{\alpha}^{\,*}_{l}$, $l=\overline{1,m}$, что задача (1.12)


будет эквивалентна детерминированной задаче


\begin{gather}


\max_{z,u}f'_{j}z=f'_{j}x_{j}^{p,\alpha}=\beta^{\alpha}_{j}(t^{p}),\notag\\


%


\Dot x =A(t)x+b(t)u,\quad x(t^p)=z, \\


%


x(t^*) \in \overline{X}_{\gamma^*}^{\,*,\alpha}=\{x \in


R^{n}:g_{*l}-\gamma+\underline{\alpha^*} \le H'_{(l)}x \le


g_l^{*}+\gamma-\overline{\alpha}\,^*_l,\ l=\overline{1,m};\ c'x \ge


\alpha-\alpha_0 \},\notag \\


%


u(t) \in U,\ \; t \in [t^p,t^*].\notag


\end{gather}





Алгоритм решения задач типа (1.13) описан в [5]. Если окажется,


что $\overline{X}_{\gamma^{*}}^{\,*,\alpha}=\emptyset$ при всех


достаточно малых $\alpha$, то увеличиваем значение $\gamma^*$.


Определим множество


$$


X^{p,\alpha}_{\gamma^p}=\{x \in R^{n}:f'_{j}x \ge


\beta^{\alpha}_{j}(t^p)-\gamma^{p},\ j =\overline{1,q}\},


$$


задав небольшое начальное значение $\gamma^{p}>0$.


Перейдем к промежутку $[t^{p-1}, t^{p}]$. Для каждого вектора из


(1.11) решим задачу


\begin{gather*}


\max_{z,u}f'_{j}z=f'_{j}x_{j}^{p-1,\alpha}=\beta^{\alpha}_{j}(t^{p-1}),\\


%


\Dot x =A(t)x+b(t)u,\quad x(t^{p-1})=z, \\


%


x(t^p) \in \overline{X}_{\gamma^p}^{\,p,\alpha}=\{x \in


R^{n}:f'_{j}x \ge


\beta^{\alpha}_{j}(t^{p})-\gamma^{p}-\alpha^{p}_{j},


\ j=\overline{1,q} \},\ \; u(t) \in U,\ \; t \in [t^{p-1},t^p],


\end{gather*}


где числа $\alpha_{j}^{p}$, $j =\overline{1,q}$, подсчитываются по


тем же правилам, что и числа $\alpha_{0}$,


$\underline{\alpha_{l}^{*}}$, $\overline{\alpha_{l}^{*}}$,


$l=\overline{1,m}$. Если


$\overline{X}_{\gamma^{p}}^{\,p,\alpha}=\emptyset$ при достаточно


малых $\alpha$, то увеличиваем $\gamma^{p}$. Продолжая этот


процесс, построим множество


$\overline{X}^{\,1,\alpha}_{\gamma^1}\ne\emptyset$. Найдем


оптимальное программное управление $\overline{u}^{0}(t)$, $t \in T$,


для состояния $x(t_*)=x_{0}$. Для этого решим сначала задачу


\begin{gather}


\gamma^{1} \rightarrow \min,\quad \Dot x=A(t)x+b(t)u,\quad x(t^{*})=x_{0},\\


%


x(t^1) \in \overline{X}_{\gamma^{1}}^{\,1,\alpha},\ \; u(t) \in U,


\ \; t \in [t_{*}, t^{1}[.\notag


\end{gather}


Если в задаче (1.14) нет допустимых управлений, то увеличиваем


$\gamma^1$. После этого находим $\overline{u}^{0}(t)$, $t \in


[t_{*}, t^{1}[$, решая задачу


\begin{gather*}


\alpha (t^1)=\max \alpha,\quad \Dot x = A(t)x+b(t)u,\quad x(t_{*})=x_{0},\\


x(t^{1}) \in \overline{X}_{\gamma^{1}}^{\,1,\alpha},\ \; u(t) \in U,


\ \; t \in [t^*, t^1[.


\end{gather*}


Пусть $\overline{x}^{0}(t^1)$ --- состояние детерминированной


системы (1.4), $x(t_*)=x_0$, в момент $t^1$ после действия


управления $\overline{u}^{0}(t)$, $t \in [t_*,t^1[$. Рассмотрим


задачу


\begin{gather}


\gamma^{2} \rightarrow \min,\quad \Dot x=A(t)x+b(t)u,


\quad x(t^{1})=\overline{x}^{0}(t^1), \\


%


x(t^1) \in \overline{X}_{\gamma^{2}}^{\,2,\alpha},\ \; u(t) \in U,


\ \; t \in [t_{1}, t^{2}[. \notag


\end{gather}


Если ограничения задачи (1.15) противоречивы, то увеличиваем


$\gamma^2$.


После этого находим управление $\overline{u}^0(t)$, $t \in


[t^1,t^2[$, решив задачу


\begin{gather*}


\alpha (t^2)=\max \alpha,\quad \Dot x = A(t)x+b(t)u,


\quad x(t_{1})=\overline{x}^{0}(t^1), \\


%


x(t^{2}) \in \overline{X}_{\gamma^{2}}^{\,2,\alpha},\ \; u(t) \in U,


\ \; t \in [t^1, t^2[.


\end{gather*}





Продолжая этот процесс, построим $\overline{u}^0(t)$, $t \in T$.


\medskip





1.2. {\it Процедура коррекции}. Поясним процедуру только для


первого промежутка. Рассмотрим промежуток $[t_*,t^1[$. Обозначим


через $J^1$ множество индексов, активных на векторе


$\overline{x}^0(t^1)$ ограничений в момент $t^1$. Построим векторы


$\Delta f^{1}_{j}=\overline{x}^0(t^1)-x^{1,\alpha}_j$, $j \in J^1 $,


и положим


\begin{equation}


\wt{f}_{j}^{1}=f_{j}+\theta\Delta f_{j}^{1},\quad j \in J^1.


\end{equation}


По векторам (1.16) найдем


\begin{equation}


\wt{\varkappa}^{0}(t^1),\ \; \wt{x}^{1,\alpha}_j,\ \; j \in


J^1 \ \;(\varkappa^{0}(t^1)=\ov{x}^{0}(t^1)).


\end{equation}


При достаточно малых $\theta > 0$


радиус $\overline{\rho}(t^1)$ сферы, в которую можно поместить


точки (1.17), будет меньше аналогичного радиуса $\rho(t^1)$ для


$\overline{x}^0(t^1)$, $x^{1,\alpha}_j$, $j \in J^1$. Выбрав точность


$\varepsilon > 0$ решения задачи (1.1), можно, продолжая коррекцию


нормалей активных ограничений, получить $\overline{\rho}(t^1) \le


\varepsilon $. Аналогичный процесс коррекции применим в остальные


моменты замыкания $t^2, \dotsc, t^p$. Управление


$\overline{u}^0(t)$, $t \in T$, при котором $\overline{\rho}(t^i)


\le \varepsilon$, $i=\overline{1,p}$, принимаем за оптимальное


программное управление $u^0(t)$, $t \in T$, задачи (1.1), если


$\gamma^* < \varepsilon$. В противном случае считаем, что в задаче


(1.1) нет допустимых управлений.





Алгоритм работы оптимального регулятора, который в режиме


реального времени строит реализацию $u^*(t)$, $t \in T$, оптимальной


обратной связи (1.9), состоит из двух частей: 1)~алгоритма


стартового этапа $(\tau=t_*)$; 2)~алгоритма текущего этапа


$\tau>t_*$. На стартовом этапе, используя программное решение


$u^0(t)$, $t \in T$, регулятор полагает $u^*(t_*)=u^0(t_*)$.


Программное решение можно получать до начала процесса управления и


поэтому нет ограничений на время его вычисления.





Пусть $\tau \in T_h$, $t ^i \le \tau < t^{i+1}$, --- текущий


момент. Если в этот момент состояние $x(\tau)$ не


доступно, то полагаем


$u^*(\tau)=u^0(\tau\mid\underline{\tau}(\tau), x(\underline{\tau}))$,


т.\,е. не производя вычислений, используем значение оптимального


программного управления, уже подсчитанного в


момент $\underline{\tau}(\tau)$ для состояния


$x(\underline{\tau})$. Рассмотрим случай


$\underline{\tau}(\tau)=\tau$, когда в момент $\tau$


доступно состояние $x(\tau)$. Это состояние


получилось из состояния $x(\underline{\tau}(\tau))$ в


результате действия на систему управления


$u^*(t)$, $t \in [\underline{\tau}, \tau[$, и возмущения


$w(t)$, $t \in [\underline{\tau}, \tau[$. При малых


$(\tau-\underline{\tau})$ состояние $x(\tau)$ мало


отличается от состояния $x_{0}(\tau)$, в которое


переходит из $x(\underline{\tau})$


детерминированная система (1.4) под действием


управления $u^*(t)$, $t \in [\underline{\tau},\tau[$. В этом


случае метод [5] быстро строит $u^0(t\mid\tau,x(\tau))$, $t \in


[\tau,t^*]$, корректируя $u^0(t\mid\underline{\tau}$,


$x(\underline{\tau}))$, $t \in [\tau,t^*]$. Если векторы


$x(\tau)$, $x_0(\tau)$ отличаются друг от друга


существенно, то время коррекции можно уменьшить


с помощью параллельных вычислений.


Используя описанный алгоритм и современные


микропроцессоры, оптимальный регулятор сможет реализовать обратную


связь для динамических систем достаточно высокого порядка.





\begin{notes}


При $T^p=T_h$ оптимальная замыкаемая обратная


связь (1.9) называется оптимальной замкнутой (``истинной'') обратной


связью. Классическим методом построения оптимальной замкнутой


обратной связи является динамическое программирование [6].


Описанный выше подход позволяет строить и оптимальные замкнутые


обратные связи. При этом требования к оперативной памяти


существенно ниже аналогичных в динамическом программировании за


счет использования процедур коррекции в процессе управления.


Чрезмерные требования к памяти в динамическом программировании


можно объяснить тем, что при его реализации дополнительная работа


в процессе управления не предусматривается, вся информация о


позиционном решении заготавливается заранее для всех мыслимых


позиций.


\end{notes}





\section{Синтез оптимальных управлений типа прямой связи}





При управлении сложными системами, функционирующими в условиях


постоянно действующих возмущений, часто доступны результаты прямых


или косвенных измерений некоторых возмущений. Использование такой


информации при формировании управляющих воздействий может


значительно повысить эффективность системы управления. Управления,


построенные по возмущениям (входным сигналам), называются


управлениями типа прямой связи или компенсационными управлениями и


отличаются от управлений типа обратной связи тем, что последние


строятся по выходным сигналам.





Заменим $w=w(t)$ в задаче (1.1) на $v=v(t)$, $t \in T$, и при


решении задачи (1.1) будем предполагать: 1)~до начала процесса


управления известно, что в процессе управления в качестве


возмущения может реализоваться любая дискретная функция $v(t) \in


V$, $t \in T$, $V=\{v \in R:|v| \le 1 \}$; 2)~в процессе управления в


каждый текущий момент времени $\tau \in T_h$ будет доступно


текущее значение $v=v(\tau) \in V $ возмущения. Погрузим задачу


(1.1) в семейство задач


\begin{gather*}


c'x(t) \rightarrow \max,\quad \Dot x=A(t)x+b(t)u+d(t)v,\ \; x(\tau)=z,


\ \; x(t^*) \in X^*,\\


u(t) \in U\ \;, v(t) \in V,\ \; t \in T(\tau)=[\tau, t^*],


\end{gather*}


зависящее от совокупности


\begin{equation}


\{\tau,z,v \},


\end{equation}


где $\tau \in T_{h}$ --- текущий момент времени, $z=x(\tau) \in R^n$


--- текущее состояние, $v=v(\tau) \in V$ --- значение возмущения


в момент $\tau$. Вектор $z$ в (2.1) считается в момент $\tau$


известным, ибо его можно построить по использованному к этому


моменту управлению $u^*(t)$, $t \in [t_*, \tau[$, и измеренному


возмущению $v^*(t)$, $t \in [t_*,z[$:


\begin{gather*}


z=F(\tau,t_*)x_0+\int_{t^*}^{\tau}F(\tau,s)b(s)u^*(s)ds+


\int_{t^*}^{\tau}F(\tau,s)d(s)v^*(s)ds \\


\big(F(t,s)=F(t)F^{-1}(s),\ \; \Dot F=A(t)F,\ \; F(0)=E\big).


\end{gather*}


Доступное управление $u_{\tau}(\cdot\mid\tau,z,v)=(u(t\mid\tau,z,v)\in


U$, $t \in T(\tau))$ назовем допустимым программным управлением для


позиции $(\tau,z,v)$, если все порожденные им, возмущением


$v(\tau)=v$ и возмущениями $v_{\tau+h}(\cdot)=(v(t) \in V$, $t \in


T(\tau+h))$ траектории


$x(t\mid\tau,z,u_{\tau}(\cdot\mid\tau,z,v),v,v_{\tau+h}(\cdot))$, $t \in


T(\tau)$, в момент $t^*$ попадают на терминальное множество $X^*$.


Обозначим


$$


X_{t^*}(u_{\tau}(\cdot\mid\tau,z,v))=\{x:x=x(t^*\mid\tau,z,v,


u_{\tau}(\cdot\mid\tau,z,v),v_{\tau+h}(\cdot)),\ v(t) \in V,\ t \in


T(\tau+h)\}.


$$


Качество допустимого управления


$u_{\tau}(\cdot\mid\tau,z,v)$ оценим числом


\begin{equation}


J(u)=\min c'x, x \in X_{t^{*}}(u_{\tau}(\cdot\mid\tau,z,v)).


\end{equation}


Допустимое управление $u_{\tau}^{0}(\cdot\mid\tau,z,v)$


считается оптимальным программным


управлением для позиции $(\tau,z,v)$, если оно доставляет


максимум величине (2.2):


$$


J(u^0)=\max_{u}J(u).


$$


Пусть $X_{\tau}$ --- множество всех $z \in R^n$, для которых


существует $u^{0}_{\tau}(\cdot\mid\tau,z,v)$ при фиксированном $\tau


\in T_{h}$ и всех $v \in V$.





\begin{defns}


Оптимальным управлением типа размыкаемой


прямой дискретной связи назовем функцию


\begin{equation}


u^{0}(\tau,z,v)=u^{0}(\tau\mid\tau,z,v),\ \; z \in X_{\tau},\ \; v \in V,


\ \; \tau \in T_{h}.


\end{equation}


\end{defns}





Траектория замкнутой системы


$$


\Dot x=A(t)x+b(t)u^{0}(t,x,v)+d(t)v,\quad x(t_*)=x_{0},


$$


определяется как решение линейного уравнения


$$


\Dot x = A(t)x+b(t)u^{*}(t)+d(t)v,\quad x(t_{*})=x_0,


$$


с управлением


\begin{gather}


u^{*}(t)=u^{0}(t_{*}+kh,x(t_*+kh), v(t_{*}+kh)),\\


t \in [t_{*}+kh,\,t_{*}+(k+1)h[,\ \; k=\overline{0,N-1}.\notag


\end{gather}


Как и в предыдущем параграфе, функцию (2.4) назовем {\it реализацией}


прямой связи (2.3). Будем говорить, что прямая связь реализуется в


режиме реального времени, если в каждый текущий момент $\tau \in


T_{h}$ время вычисления значения $u^{*}(\tau)$ по известным $z,\ v$


не превосходит $h$. Устройство, способное выполнить эту работу,


называется оптимальным регулятором.





Определим множества


\begin{gather*}


X_{\tau}^{1}=\bigg\{x \in R^n:x=\int_{\tau+h}^{\tau}F(t^*,s)d(s)v(s)ds,


v(t) \in V,\ t \in T(\tau+h)\bigg\}, \\


%


X_{\tau}^{2}=\{x \in R^{n}: x+X^1_{\tau} \subset X^*\}, \\


%


X_{\tau}^{v,z}=\bigg \{x \in R^n:x=F(t^*,\tau)z+


v\int_{\tau}^{\tau+h}F(t^{*},s)d(s)ds,\ v \in V\bigg\}, \\


%


X^{3}_{\tau}=\bigg\{x \in R^{n}:x=\int_{\tau}^{t^*}F(t^*,s)b(s)u(s)ds,


u(t) \in U,\ t \in T(\tau)\bigg\}.


\end{gather*}





В терминах этих множеств условие допустимости управления


$u_{\tau}(\cdot\mid\tau,z,v)$ принимает вид $X_{\tau}^{v,z} \subset


(X_{\tau}^{2}-X_{\tau}^3)$. Аналогичным образом определим


множество $X_{\tau}^{2,\alpha}$, заменив $X^{*}$ на


$X^{*,\alpha}=X^* \cap \{x \in R^n:c'x \ge \alpha\}$.


Максимальное $\alpha=\alpha^0$, при котором выполняется включение


\begin{equation}


X_{\tau}^{v,z} \subset (X_{\tau}^{2,\alpha}-X_{\tau}^{3}),


\end{equation}


равно оптимальному гарантированному значению критерия качества


задачи. Управление, с которым соотношение (2.5) выполняется при


$\alpha=\alpha^0$, равно $u^0(t\mid\tau,z,v)$, $t \in T(\tau)$.





\begin{notes}


Нетрудно заметить, что определения и метод


решения принципиально не изменятся, если предположение~2) заменить


на условие, что в процессе управления в каждый текущий момент $\tau


\in T_{h}$ по результатам измерения до момента $\tau$ включительно


доступен прогноз $\wh{v}(t)$, $t \in [\tau,\tau+\theta[$,


$\theta=lh$, $l=\overline{1,l^*}$, возмущения $v(t)$, $t \in [\tau,


\tau+\theta]$, с ошибкой $\overline{v}(t)$, $t \in [\tau,


\tau+\theta]:$ $v(t)=\wh{v}(t)+\wt{v}(t)$, $|\wh{v}(t)| \le


\ov{v}(t)$, $t \in [\tau, \tau+\theta]$.


\end{notes}





Для конструктивного построения оптимального программного


управления и на его базе оптимального управления типа прямой связи


используем полиэдральную аппроксимацию соотношения (2.5), решаем


аппроксимационную задачу и доводим полученное решение до решения


исходной задачи (1.1) (см. \S\,1 и предположения \S\,2). Как отмечено


выше, реализация оптимальных управлений типа размыкаемой связи


основана на быстрых коррекциях программных решений. Информация о


некоторых элементах программного решения, полученная в позиции


$(\tau,x^{*}(\tau),v^{*}(\tau))$, преобразуется двойственным


методом [5] в информацию для позиции


$(\tau+h,x^*(\tau+h),v^*(\tau+h))$, когда процесс управления


перешел от момента $\tau$ в следующий момент $\tau+h$ и стали


известны $x^{*}(\tau+h)$, $v^{*}(\tau+h)$. Работу по коррекции


информации современные микропроцессоры могут выполнить за время,


меньшее чем $h$, для систем управления весьма высокого порядка.


Это позволяет построить на их основе алгоритм функционирования


оптимального регулятора.





\begin{notes}


Если дополнить априорную информацию по


рассматриваемой задаче и считать, что до начала процесса


управления в моменты замыкания $T^p=\{t^i \in T_{h},


\ i=\overline{1,p}\}$, $t_{*}<t^1< \dotsb < t^p < t^*$, будут известны


значения $v(t^i)$, $i=\overline{1,p}$, возмущения, то можно ввести


понятие оптимальной замыкаемой прямой дискретной связи.


Использование таких управлений расширяет множество начальных


состояний $x_{0}$, для которых существует решение задачи


гарантированной оптимизации. Алгоритм работы оптимального


регулятора, реализующего оптимальное управление типа замыкаемой


прямой связи, строится по схеме алгоритма из \S\,1.


\end{notes}





\section{Синтез оптимальных управлений типа комбинированной связи}





3.1. {\it Метод решения.} При управлении динамическими


системами, функционирующими в условиях


постоянно действующих возмущений, кроме


принципа обратной связи и принципа прямой связи


используется и комбинированный принцип


(принцип прямо-обратной связи). Как отмечалось выше,


комбинированный принцип управления базируется на обоих первых


принципах и реализуется в виде прямых и обратных связей,


использующих как входные, так и выходные сигналы. Ниже результаты


\S\S\,1,~2 развиваются на случай оптимального комбинированного


управления.





Пусть в классе дискретных управлений $u(t)$, $t \in T$, возмущений


$v(t)$, $t \in T$, и кусочно-непрерывных возмущений $w(t)$, $t \in T$,


рассматривается задача


\begin{gather}


c'x(t^*)\rightarrow \max,\quad \Dot x = A(t) x+b(t)u+d(t)v+f(t)w,


\ \; x(t_{*})=x_{0}, \\


%


x(t^*) \in X^*=\{x \in R^n:g_{*} \le Hx\le g^*\},\quad u(t) \in U=\{u


\in R:|u| \le 1 \}, \notag \\


%


v(t) \in V=\{v \in R:|v| \le 1\},\quad w(t) \in W=\{w \in R:|w|\le 1,


\ t \in T \} \notag


\end{gather}


($x \in R^n$, $H \in R^{m \times n}$, $A(t)$, $b(t)$, $d(t)$,


$f(t)$, $t \in T$, --- кусочно-непрерывные функции).





Предположим, что до начала процесса управления известно, что 1)~в


процессе управления могут реализоваться любые дискретные


возмущения $v(t) \in V$, $t \in T$, и кусочно-непрерывные возмущения


$w(t) \in W$, $t \in T$; 2)~в каждый текущий момент времени $\tau


\in T_{h}$ будут известны текущие состояние $x(\tau)$ и


значение $v(\tau)$ возмущения $v(t)$, $t \in T$. Погрузим задачу


(3.1) в семейство аналогичных задач


\begin{gather}


c'x(t^*)\rightarrow \max,\quad \Dot x = A(t) x+b(t)u+d(t)v+f(t)w,


\ \; x(\tau)=z,\ \; v(\tau)=v, \\


%


x(t^*) \in X^*,\ \; u(t) \in V,\ \; v(t)\in U,


\ \; w(t) \in U,\ \; t \in T(\tau+h),\notag


\end{gather}


зависящее от тройки $\tau \in T_{h}$, $z \in R^n$, $v \in V$.





Допустимое управление $u(\cdot\mid\tau,z,v)=(u(t\mid\tau,z,v)


\in U$, $t \in T(\tau))$ назовем допустимым


программным управлением для позиции $(\tau,z,v)$,


если траектории $x(t\mid\tau,z,v,\ u(\cdot\mid\tau,z,v),


\ v_{\tau+h}(\cdot),\ w_{\tau}(\cdot))$, $t \in T(\tau)$, системы


(3.2), порожденные этим управлением и всеми


возможными возмущениями $v_{\tau+h}(\cdot)=(v(t),\ t


\in T(\tau+h))$, $w_{\tau}(\cdot)=(w(t)$, $t \in T(\tau))$, в


момент $t^*$ попадают на терминальное множество


$X^*$. Обозначим $X_{t^*}(u(\cdot\mid\tau,z,v))=\{x \in


R^n:x=x(t^*\mid\tau,z,v,\ u(\cdot\mid\tau,z,v),


\ v_{\tau+h}(\cdot),\ w_{\tau}(\cdot))$, $v(t) \in V$, $t \in T(\tau+h)$;


$w(t) \in W$, $t \in T(\tau)\}$. Качество допустимого


управления оценим числом


\begin{equation}


J(u(\cdot\mid\tau,z,v))=\min c'x,\quad x \in X_{t^*}(u(\cdot\mid\tau,z,v)).


\end{equation}





Допустимое управление $u^0(t\mid\tau,z,v)$, $t \in


T(\tau)$, будем называть оптимальным программным


управлением задачи (3.1), если оно среди всех


допустимых управлений доставляет наибольшее


значение величине (3.3):


$J(u^0(\cdot\mid\tau,z,v))=\max\limits_{u}J(u(\cdot\mid\tau,z,v))$.





Пусть $X_{\tau}$ --- множество векторов $z \in R^n$, для которых


задача (3.2) имеет программное решение $u^0(t\mid\tau,z,v)$, $t \in


T(\tau)$, при каждом фиксированном $\tau \in T_h$ и всех $v \in V$.





\begin{defns}


Функция


\begin{equation}


u^0(\tau,z,v)=u^0(\tau\mid\tau,z,v),\ \; z \in X_{\tau},\ \; v \in V,


\ \;\tau \in T_{h},


\end{equation}


называется оптимальным управлением типа


прямо-обратной (комбинированной) дискретной


связи для задачи (3.1).


\end{defns}





Траекторию замкнутой нелинейной системы $\Dot


x=A(t)x+b(t)u^0(t,x,v)+d(t)v+f(t)w$, $x(t_{*})=x_{0}$, которая


получается из (3.1) после подстановки в нее вместо управления $u$


функции (3.4), определим как решение линейного уравнения $\Dot x


=A(t)x+b(t)u^*(t)+d(t)v+f(t)w$, $x(t_{*})=x_{0}$, с управлением


\begin{gather*}


u^*(t)=u^0(t_{*}+kh,\ x(t_*+kh),\ v(t_*+kh)),\\


t \in [t_{*}+kh,t_{*}+(k+1)h[,\quad k =\overline{0,N-1}.


\end{gather*}





Аналогично изложенному выше вводится понятие {\it реализации}


прямо-обратной связи (3.4) в конкретном процессе управления.





Введем множества


\begin{gather}


X_{\tau}^{1}=\bigg\{x \in R^n: x =


\int_{\tau-h}^{t^*}F(t^*,s)d(s)v(s)ds+


\int_{\tau}^{t^*}F(t^*,s)f(s)w(s)ds, \notag \\


%


v(t) \in V,\ \; t \in T(\tau+h),\ \; w(t) \in W,


\ \;t \in T(\tau) \bigg\},\notag \\


%


X_{\tau}^{2}=\{x \in R^n:x+X_{\tau}^1 \subset X^*\},\notag \\


%


X^{v,z}_{\tau}=\bigg\{x \in


R^n:x=F(t^*,\tau)z+v\int_{\tau}^{\tau+h}F(t^*,s)d(s)ds,\ v \in V


\bigg\}, \\


%


X_{\tau}^3=\bigg\{x \in R^n:x=\int_{\tau}^{t^*}F(t^*,s)b(s)u(s)ds,


\ u(t) \in U,\ t \in T(\tau) \bigg\}.\notag


\end{gather}





В терминах множеств (3.5) условие допустимости программного


управления принимает вид $X_{\tau}^{v,z} \subset


(X_{\tau}^2-X_{\tau}^3)$. Геометрическая интерпретация условия


оптимальности допустимого управления получается следующим образом.


Введем множество $X^{*,\alpha}=X^{*}\cap\{x \in R^n:c'x \ge


\alpha\}$. Положим вместо $X_{\tau}^2$: $X_{\tau}^{2,\alpha}=\{x


\in R^n: x + X_{\tau}^{1} \subset X^{*,\alpha}\}$. Допустимое


программное управление обеспечивает критерию качества задачи (3.1)


гарантированное значение $\alpha$, если


\begin{equation}


X_{\tau}^{v,z} \subset (X_{\tau}^{2,\alpha}-X_{\tau}^{3}).


\end{equation}


Наибольшее $\alpha=\alpha^{0}$, при котором выполняется включение


(3.6), равно оптимальному (наибольшему) гарантированному значению


критерия качества. Допустимое управление, на котором выполняется


условие (3.6) с $\alpha=\alpha_{0}$, и есть оптимальное


программное управление для позиции $(\tau,z,v)$.





Конструктивная проверка соотношения (3.6) и построение


оптимального программного управления состоят из двух процедур:


1)~решение вспомогательных линейных задач оптимального управления,


построенных на базе полиэдральных аппроксимаций множеств из (3.5),


и замене соотношения (3.6) на приближенное


$$


\max_{x \in X_{\tau}^{v,z}}f'_{j}x \le \max_{x \in


X_{\tau}^{2,\alpha},\ y \in X_{\tau}^{3}}f'_{j}(x-y),\quad


j=\overline{1,q},


$$


где $f_{j}$, $j=\overline{1,q}$, --- набор единичных


векторов (см. \S\,1); 2)~коррекции (доводки),


представляющей процесс последовательного


улучшения решений вспомогательных задач до


построения решения задачи (3.1) с заданной


точностью.





Алгоритм работы оптимального регулятора представляет процедуру 


коррекции в текущей позиции $(\tau,x^*(\tau),v^*(\tau))$


элементов программного решения, построенного для предыдущей


позиции $(\tau-h,x^*(\tau-h),v^*(\tau-h))$. В построении


оптимальных программных управлений, реализующих оптимальные


управления типа комбинированной связи (3.4), определяющую роль


играют быстрые алгоритмы [5].


\medskip





3.2. {\it Замкнутая комбинированная связь и


метод динамического программирования.}





\begin{defns}


При $T^p=T_h$ оптимальное управление типа


замыкаемой комбинированной связи назовем оптимальным управлением


типа замкнутой комбинированной связи.


\end{defns}





Замкнутые связи используют наибольший объем априорной информации и


поэтому считаются наиболее эффективными. Поскольку выше, при


изложении метода синтеза оптимальных замыкаемых комбинированных


связей, условие $T^p\ne T_h$ нигде не использовалось, то


синтез оптимальных управлений типа замкнутых комбинированных


связей производится по описанной выше схеме.





Классическим методом синтеза оптимальных замкнутых связей является


динамическое программирование [6]. Для исследуемой задачи этот


метод приводит к уравнению Беллмана


\begin{multline}


B_{\tau}(z)=\min_{v \in V}\ \max_{u \in U}\; \min_{w(s) \in W,\ s \in


[\tau, \tau+h]}


B_{\tau+h}\bigg(F(\tau+h,\tau)z+u\int_{\tau}^{\tau+h}F(\tau+h,s)b(s)ds+ \\


%


+v\int_{\tau}^{\tau+h}F(\tau+h,s)d(s)ds+


\int_{\tau}^{\tau+h}F(\tau+h,s)f(s)w(s)ds\bigg)


\end{multline}


с начальным условием


$$


B_{t^*}(z)=


\begin{cases}


c'z, & z \in X^*; \\ -\infty, & z \notin X^*,


\end{cases}


$$


и к соотношению для вычисления оптимального управления типа


замкнутой комбинированной связи $u^0(\tau,z,v)$, $z \in X_{\tau}$,


$v \in V$, $\tau \in T_{h}$,


\begin{multline}


\min_{w(s) \in W,\ s \in [\tau, \tau+h]}


B_{\tau+h}\bigg(F(\tau+h,\tau)z+u^0(\tau,z,v)


\int_{\tau}^{\tau+h}F(\tau+h,s)b(s)ds+ \\


%


+v\int_{\tau}^{\tau+h}F(\tau+h,s)d(s)ds+


\int_{\tau}^{\tau+h}F(\tau+h,s)f(s)w(s)ds\bigg)= \\


%


=\max_{u \in U}\;\min_{w(s) \in W,\ s \in [\tau, \tau+h]}


B_{\tau+h}\bigg(F(\tau+h,\tau)z+ u


\int_{\tau}^{\tau+h}F(\tau+h,s)b(s)ds+ \\


%


+v\int_{\tau}^{\tau+h}F(\tau+h,s)d(s)ds+


\int_{\tau}^{\tau+h}F(\tau+h,s)f(s)w(s)ds\bigg).


\end{multline}


Здесь $F(t,s)=F(t)F^{-1}(s)$, $\Dot F=A(t)F$, $F(0)=E$, $X_{\tau}=\{z \in


R^n:B_{\tau}(z)\ne -\infty \}$. Как видно из (3.7), (3.8),


уравнение Беллмана можно решить по шагам справа


налево, табулируя попутно функцию Беллмана


$B_{1}(z)$, $z \in X_{\tau}$, $\tau \in T_{h}$, и оптимальное


управление $u^0(\tau,z,v)$, $z \in X_{\tau}$, $v \in V$, $\tau \in


T_{h}$. Эту работу можно проделать до начала процесса управления,


поскольку в (3.7), (3.8) используется только априорная


информация. Однако, как хорошо известно из [6],


реализовать метод для построения высокоточных решений задачи c $n


\ge 3$ чрезвычайно трудно из-за чрезмерного объема требуемой


оперативной памяти ЭВМ.





Принципиальное отличие метода, изложенного в данной работе, от


метода динамического программирования состоит в том, что на


предварительном этапе (до начала процесса управления)


заготавливается грубая информация о решении с использованием


небольшого объема оперативной памяти, а потом, в процессе


управления, эта информация в каждый текущий момент времени $\tau


\in T_{h}$ быстро корректируется [5] в зависимости от


реализующихся состояния $x^*(\tau)$ и возмущения $v^*(\tau)$.


B [5], [7] разработан ряд быстрых алгоритмов,


которые позволяют


осуществить коррекцию в режиме реального времени на современных


вычислительных устройствах для систем управления достаточно


высокого порядка.








\begin{thebibliography}{9}





\bibitem{1}


Фельдбаум А.А. {\it Основы теории оптимальных автоматических


систем}. -- М.: ГИФМЛ, 1963. -- 294 с.


\bibitem{2}


{\it Фильтрация и стохастическое управление в динамических


системах} / Под ред. Ле\-он\-де\-са~К.Т. -- М.: Мир, 1980. -- 404 с.


\bibitem{3}


Габасов Р., Кириллова Ф.М., Костина Е.А. {\it Замыкаемые


обратные связи по состоянию для оптимизации неопределенных систем


управления} // Автоматика и телемеханика. -- 1996.-- \No\,7. --


С.\,121--130.


\bibitem{4}


Габасов Р., Кириллова Ф.М., Костюкова О.И. {\it Алгоритм


оптимизации в режиме реального времени не полностью определенной


линейной системы управления} // Автоматика и телемеханика. --


1993. -- \No\,4. -- С.\,34--43.


\bibitem{5}


Балашевич Н.В., Габасов Р., Кириллова Ф.М. {\it Численные


методы программной и позиционной оптимизации линейных систем


управления}


// Журн. вычисл. матем. и матем. физ. -- 2000. --


Т.\,40. -- \No\,6. -- С.\,799--819.


\bibitem{6}


Беллман Р. {\it Динамическое программирование.} -- М.: Ин. лит.,


1960. -- 400 с.


\bibitem{7}


Gabasov R., Kirillova F.M., Balashevich N.V. {\it On the


synthesis problem for optimal control systems} // SIAM J.


Control and Optim. -- 2001. -- V.\,39. -- P.1008--1042.





\end{thebibliography}





\vskip 2cm





\post{\it Белорусский государственный университет}{\it Поступила}


\post{\it Институт математики}{15.08.2001}


\post{\it Национальной Академии наук Беларуси}{}





\label{end}


\end{document}





y


